\documentclass[10pt]{article}
\lineskip=0pt

% Formatting follows st2334-midterms-cheatsheet.tex.
\usepackage[skins,breakable,raster]{tcolorbox}
\tcbsetforeverylayer{
  skin=enhanced jigsaw,
  fonttitle=\bfseries\small,
  parbox=false,
  colback=white,
  subtitle style={colback=gray!15,fontupper=\scriptsize\bfseries},
  left=1pt,right=1pt,top=1pt,bottom=1pt,
  boxrule=0.35pt,
  nobeforeafter,
  breakable, height fixed for=first,
  pad at break*=0pt, pad before break*=0pt, pad after break=0pt,
  break at=-\baselineskip/0pt
}
\usepackage{parskip}
\setlength\parskip{0pt}
\usepackage{multicol}
\usepackage{geometry}
\newgeometry{noheadfoot=true,tmargin=5pt,bmargin=5pt,hmargin=5pt}
\geometry{paperheight=297mm,paperwidth=420mm}
\usepackage{xcolor}
\usepackage{soul}
\usepackage[fleqn]{amsmath}
\usepackage{amssymb,mathtools}
\usepackage[inline]{enumitem}
\usepackage{tabularx}
\usepackage{tabularray}
\usepackage{graphicx}
\pagestyle{empty}
\setlength{\parindent}{0pt}
\setlength{\columnsep}{2pt}
\setlist{leftmargin=*,labelsep=2pt,itemsep=0pt,topsep=0pt,parsep=0pt}
\setlength{\abovedisplayskip}{1pt}
\setlength{\belowdisplayskip}{1pt}
\setlength{\abovedisplayshortskip}{0pt}
\setlength{\belowdisplayshortskip}{0pt}
\newcommand{\warn}[1]{\textbf{\color{red!65!black}#1}}
\newcommand{\key}[1]{\textbf{#1}}
\newcommand{\ind}[1]{\mathbf 1_{\{#1\}}}
\newcommand{\boxstart}[1]{\begin{tcolorbox}[title={#1},sharp corners]}
\newcommand{\boxend}{\end{tcolorbox}}

% Sources: lecture notes Ch. 1--7; slides Weeks 1--6 and 8--12;
% Week 8 quiz; Tutorials 8--9. FormulaSheet_Tables.pdf is the companion.
% Supplied PMFs/PDFs, moments, CI/test tables and numerical tables are omitted.
% Short substitutions below illustrate methods rather than repeat those tables.
\begin{document}
\small
% Font-size commands reset display spacing; apply the template's compact
% spacing after \small so it takes effect at the unchanged 9pt body size.
\setlength{\abovedisplayskip}{1pt}
\setlength{\belowdisplayskip}{1pt}
\setlength{\abovedisplayshortskip}{0pt}
\setlength{\belowdisplayshortskip}{0pt}

% PAGE 1 / COLUMN 1: counting and conditional probability.
\begin{multicols*}{3}
\vspace*{-1em}
\boxstart{ST2334 Finals --- Ch. 1--4: Probability}
\textbf{Companion: FormulaSheet\_Tables.pdf (FS).} FS p.\,1 supplies basic probability, moments, distributions and approximations; p.\,2 supplies mean CI/test formulas; subsequent pages contain tables. Here: methods, missing results and exam traps. Chapter numbers follow the notes, not lecture weeks.
\boxend

\boxstart{Ch. 1 --- Counting: set up before calculating}
\begin{enumerate}
\item Define equally likely outcomes; decide whether order matters and whether replacement is allowed.
\item Split into disjoint cases. For ``at least one'', count the complement; for restrictions, count all minus forbidden.
\item Choose groups/values, then positions. Divide only for ordering counted more than once; check the total against the sample space.
\end{enumerate}
\textbf{Five-digit strings (leading zero allowed):} repeated-digit patterns give the following counts out of $10^5$:
\[
\begin{array}{c|r@{\qquad}c|r}
1+1+1+1+1&30240&2+1+1+1&50400\\
2+2+1&10800&3+1+1&7200\\
3+2&900&4+1&450\\
5&10&&
\end{array}
\]
E.g. two pairs: $\binom{10}{2}\,8\,5!/(2!2!)=10800$. The two pair values were already chosen as an unordered set: no extra $2!$ division. For alternating $n$ distinct people of each of two types in a line: $2(n!)^2$.
\boxend

\boxstart{Ch. 1 --- Collisions and identifier length}
For $k$ iid assignments into $m$ equally likely slots, $k\le m$,
\[
P(\text{collision})=1-\prod_{i=0}^{k-1}(1-i/m)
\approx1-e^{-k(k-1)/(2m)}.
\]
For $k>m$, collision is certain. Any-pair collision needs order $\sqrt m$ draws; hitting one fixed target needs order $m$.
\textbf{Guaranteed bound:} $P(\text{collision})\le\binom{k}{2}/m$; for alphabet size $b$ and length $\ell$, choose $b^\ell\ge\binom{k}{2}/\varepsilon$ to bound risk by $\varepsilon$. Round $\ell$ up. The exponential approximation is useful when $k/m$ is small.
\boxend

\boxstart{Ch. 1 --- Uniform shuffle and fixed points}
\textbf{Fisher--Yates:} at $i=n-1,\ldots,1$, choose $j$ uniformly from $0,\ldots,i$ and swap. The $n!$ equally likely choice sequences correspond to the $n!$ permutations. Choosing from the full range at each step is generally biased.

For $A_i=$ ``item $i$ stays fixed'', an intersection of $r$ specified events has probability $(n-r)!/n!$. Inclusion--exclusion gives
\[
P(\text{some fixed point})=\sum_{r=1}^n\frac{(-1)^{r+1}}{r!}
\longrightarrow1-e^{-1}.
\]
A correct shuffle need not move every item. General inclusion--exclusion alternates sums of single events, pair intersections, triples, etc.
\boxend

\boxstart{Ch. 1 --- Bayes, paths and repeated evidence}
\begin{enumerate}
\item List mutually exclusive causes $H_i$ and their priors.
\item Compute each weight $w_i=P(H_i)P(E\mid H_i)$.
\item Normalise: posterior for $H_i$ is $w_i/\sum_jw_j$.
\end{enumerate}
\textbf{Several observations:} multiply their likelihoods only if they are conditionally independent \emph{given the cause}. In naive Bayes, use a likelihood for each present \emph{and absent} feature within each class. Sensitivity is not the posterior probability of disease.
\[
\text{posterior odds}=\text{prior odds}\times\text{likelihood ratio}.
\]
For a tree, multiply along paths and add the paths consistent with the evidence. Mixing different causes can create marginal dependence even when observations are conditionally independent.
\boxend

\boxstart{Ch. 1 --- Independence and fair-bit extraction}
\textbf{Pairwise} checks each pair; \textbf{mutual} checks every subcollection. If independent fair bits $X,K$ give $C=X\oplus K$, every pair is independent, but the triple is not. Reusing $K$ leaks $C_1\oplus C_2=X_1\oplus X_2$.

\textbf{Fair extractor:} iid biased bits, $P(1)=p$: map $01\to0$, $10\to1$, discard $00,11$. Both accepted pairs have probability $p(1-p)$; acceptance probability is $2p(1-p)$. For $0<p<1$, expected input bits per output bit is $1/[p(1-p)]$. Unequal biases or dependence can invalidate fairness.
\boxend

\boxstart{Ch. 1/5 --- Fixed-size sampling without replacement}
\textbf{Algorithm S:} with $t$ records remaining and $m$ still needed, keep the next with probability $m/t$; decrement $m$ if kept and always decrement $t$. Every complete $n$-subset has the same probability.

\textbf{Hypergeometric counting:} among $N$ objects, $K$ marked; choose $n$ without replacement. The count $X$ of marked objects has
\[
P(X=x)=\frac{\binom Kx\binom{N-K}{n-x}}{\binom Nn},\quad
\max(0,n-N+K)\le x\le\min(n,K).
\]
Use $E(X)=nK/N$. Do not model these dependent draws as binomial trials. Example: $N=18,K=5,n=4$ gives $P(X=2)=13/51$.
\boxend

\columnbreak
% PAGE 1 / COLUMN 2: random variables and dependence.
\boxstart{Ch. 2 --- PMF/PDF/CDF question workflow}
\begin{enumerate}
\item Write the support; normalise there to find any constant.
\item Find the CDF below/inside/above support, using only allowed values.
\item Discrete endpoints matter; continuous endpoints have probability zero.
\item For a quantile, solve the CDF condition; for discrete $X$, take the smallest $x$ with $F(x)\ge p$.
\end{enumerate}
Density height may exceed 1. Finite mean need not imply finite variance: $f(x)=2/x^3$, $x\ge1$, has mean 2 but infinite second moment.
\boxend

\boxstart{Ch. 2 --- Transformations, extrema and simulation}
\textbf{CDF method:} find the new support, express $\{g(X)\le y\}$ in terms of $X$, then calculate its probability. Differentiate if needed. Non-monotone maps may need several intervals.
For a differentiable one-to-one map,
\[
f_Y(y)=f_X(g^{-1}(y))\left|\frac{d}{dy}g^{-1}(y)\right|.
\]
For iid $X_1,\ldots,X_n$ with CDF $F$,
\[
P(\max X_i\le x)=F(x)^n,\quad
P(\min X_i>x)=[1-F(x)]^n.
\]
\textbf{Inverse CDF:} draw $U\sim U(0,1)$ and set $X=F^{-1}(U)$. Example: $F(x)=1-x^{-2}$ on $x\ge1$ gives $X=(1-U)^{-1/2}$. Verify by computing the generated CDF. For jumps use $F^{-1}(u)=\inf\{x:F(x)\ge u\}$.
\boxend

\boxstart{Ch. 2 --- Expectations: choose an easy representation}
\textbf{Transformed quantity:} evaluate $E[g(X)]$ directly using FS; generally $E[g(X)]\ne g(E[X])$. For quadratic cost, use $E(X^2)=V(X)+[E(X)]^2$ before optimising.

\textbf{Indicators:} write a count as $N=\sum I_i$; sum single-event probabilities for its mean. For variance, add all individual variances and twice every pair covariance; compute each covariance from $P(A_i\cap A_j)-P(A_i)P(A_j)$.

\textbf{Tail method:} use the integer tail-sum in FS when $P(X\ge k)$ is easier than the PMF. For continuous nonnegative $X$,
\[
E(X)=\int_0^\infty P(X>t)\,dt,\quad
E[\min(X,T)]=\int_0^T P(X>t)\,dt.
\]
For a uniformly chosen $m$-bit counter with wraparound, $P(\text{flips}\ge k)=2^{-(k-1)}$, $1\le k\le m$, so the mean is $2-2^{1-m}$ (check $m=1$).
\boxend

\boxstart{Ch. 2 --- Records and coupon collection}
\textbf{Records in a uniform random permutation:} record indicator $I_i$ has probability $1/i$. Pairwise independence yields
\[
E(R)=H_n,\qquad V(R)=\sum_{i=1}^n(1/i-1/i^2),\quad H_n=\sum_{i=1}^n1/i.
\]
\textbf{Coupon collector:} with $d$ equally likely designs, wait for the next new type after $j$ collected: success probability $(d-j)/d$. Sum the geometric means to obtain $E(T)=dH_d$. Unequal design probabilities invalidate this shortcut.

\boxend

\boxstart{Ch. 3 --- Joint support, marginals and conditionals}
Draw the region first. Normalise; integrate/sum out the unwanted coordinate for a marginal. Divide by the conditioning marginal where positive. Compute conditional means using the conditional density.

\textbf{Triangle example:} $f(x,y)=c$ on $0<x<y<1$ gives $c=2$, $f_X(x)=2(1-x)$ and $f_Y(y)=2y$. Thus $X\mid Y=y$ is uniform on $(0,y)$ and $E(X\mid Y=y)=y/2$.

An event such as $X+Y<a$ adds a boundary; split integrals where limits change. Independence requires factorisation including the support: a constant density on a triangle does \emph{not} imply independence.
\boxend

\boxstart{Ch. 3 --- Dependence patterns and checks}
Zero covariance need not mean independence: symmetric $X$ and $Y=X^2$ provide a counterexample (assuming the required moments exist).

\textbf{Fixed-size sample:} if exactly $n$ of $N$ records are chosen, for distinct inclusion indicators,
\[
P(I_i=I_j=1)=\frac{n(n-1)}{N(N-1)},\quad
\operatorname{Cov}(I_i,I_j)=-\frac{n(N-n)}{N^2(N-1)}.
\]
Their sum is fixed, so its variance is 0: a useful covariance check.

\textbf{Runs in a random permutation:} $R=1+\sum_{i=1}^{n-1}\ind{A_i>A_{i+1}}$. Each descent has probability $1/2$; adjacent descents have joint probability $1/6$ and covariance $-1/12$; non-adjacent ones are independent. For $n\ge2$,
\[
E(R)=\frac{n+1}{2},\qquad V(R)=\frac{n+1}{12}.
\]
\boxend

\boxstart{Ch. 3/4 --- Independent Poissons given their sum}
If $X\sim\mathrm{Pois}(\lambda)$ and $Y\sim\mathrm{Pois}(\mu)$ are independent,
\[
X\mid(X+Y=n)\sim\mathrm{Bin}\left(n,\frac{\lambda}{\lambda+\mu}\right).
\]
Derive by $P(X=x,Y=n-x)/P(X+Y=n)$, $0\le x\le n$. Conditioning on the total makes the components dependent. For rates 2 and 3 and total 10, the conditional mean of the first is $10(2/5)=4$.
\boxend

\columnbreak
% PAGE 1 / COLUMN 3: special models and lecture applications.
\boxstart{Ch. 4 --- Choose the model from the experiment}
\begin{tabularx}{\linewidth}{@{}lX@{}}
Bernoulli / binomial & One trial / successes in fixed $n$ iid trials.\\
Geometric / neg. bin. & Trials \emph{including} the first / $k$th success.\\
Poisson & Count over fixed time/space; scale rate to that window.\\
Poisson-binomial & Independent Bernoulli sum with unequal $p_i$.\\
Hypergeometric & Marked count sampled without replacement.\\
Uniform & Constant density over an interval; length ratios.\\
Exponential & Waiting time at a constant rate.\\
Normal & Symmetric continuous measurements; variance is the second parameter.
\end{tabularx}
\textbf{Check:} fixed count or stopping time? independence? equal success probabilities? For ``$k$th success on trial $r$'', the first $r-1$ trials contain exactly $k-1$ successes and trial $r$ succeeds. Convert failures-only counts to the FS trials-including-success convention.
\boxend

\boxstart{Ch. 4 --- Approximations and tail wording}
\textbf{Normal approximation:} both expected successes and failures must be sufficiently large; large $n$ alone is insufficient. \textbf{Poisson approximation:} individual success probabilities should be small; sum them for the parameter (see FS).

For integer-valued $X$, first convert strict bounds: $X>k$ means $X\ge k+1$; $X<k$ means $X\le k-1$. Then use the FS continuity correction \emph{once}, before standardising. Never continuity-correct a genuinely continuous measurement.

Standardise with the \emph{SD}, not variance; $P(|Z|>c)$ has two tails. For inverse questions, convert to the table's tail convention before solving for the threshold.
\boxend

\boxstart{Ch. 4 --- Survival and remaining lifetime}
For $s$ with $P(X>s)>0$, use the remaining support:
\[
P(X>s+t\mid X>s)=\frac{P(X>s+t)}{P(X>s)},\quad t\ge0.
\]
\textbf{Uniform $U(0,b)$:} given $X>s$ with $0\le s<b$, residual time is $U(0,b-s)$ with mean $(b-s)/2$.
\textbf{Exponential:} residual time has the original distribution.
\textbf{Pareto $f(x)=2/x^3$, $x\ge1$:} for $s\ge1$, conditional mean is $2s$, so mean remaining time is $s$. For $0\le s<1$, it is $2-s$.

\textbf{One fresh attempt versus continuing:} compare $r(s)=E(X-s\mid X>s)$ with $E(X)$, assuming no restart cost and an independent fresh draw. In this Pareto example, restarting once is better for $s>2$.
\boxend

\boxstart{Ch. 4 --- Repeated timeouts: renewal equation}
If every attempt is stopped at $T$ and restarted independently without cost, with $P(X\le T)>0$,
\[
m(T)=E[\min(X,T)]+P(X>T)m(T)
\quad\Rightarrow\quad m(T)=\frac{E[\min(X,T)]}{P(X\le T)}.
\]
Exponential timeouts leave the mean unchanged. For the Pareto above, $T>1$,
\[
m(T)=\frac{2-1/T}{1-1/T^2},\qquad T_{\rm opt}=2+\sqrt3\approx3.732.
\]
Derive by differentiating and checking the domain/boundaries. The optimal repeated-timeout policy is a different question from comparing the residual mean with one fresh attempt.
\boxend

\boxstart{Ch. 4 --- Packet retransmission and optimisation}
For independent bit-error probability $p$, a packet of $\ell$ bits succeeds with $s=(1-p)^\ell$. Attempts through success have geometric mean $1/s$; through $k$ successes, negative-binomial mean $k/s$.

With $h$ header bits and payload $\ell-h$, expected transmitted bits per useful bit are
\[
C(\ell)=\frac{\ell}{(\ell-h)(1-p)^\ell},\qquad \ell>h.
\]
Let $a=-\ln(1-p)>0$. Solve $(\log C)'=1/\ell-1/(\ell-h)+a=0$: $\ell_*=(h+\sqrt{h^2+4h/a})/2$. Check feasible neighbouring integers and constraints. For $p\ll1$, $a\approx p$.
\boxend

\boxstart{Ch. 4/5 --- Reservoir sampling and replacement counts}
Keep a uniform reservoir of $k$ records. On arrival $i>k$, draw $J_i$ uniformly in $1,\ldots,i$; replace slot $J_i$ iff $J_i\le k$. Each record is included with probability $k/i$ after $i$ arrivals.

For arrivals $m+1,\ldots,n$, $m\ge k$, replacement decisions are independent Bernoulli$(k/i)$. Thus
\[
E(R)=k\sum_{i=m+1}^n\frac1i,\qquad
V(R)=E(R)-k^2\sum_{i=m+1}^n\frac1{i^2}.
\]
For large $m$, $E(R)\approx k\log(n/m)$; during a doubling, approximately $k\log2$. Apply the FS Poisson-binomial bound if using a Poisson approximation. Reservoir contents themselves are dependent.
\boxend

\boxstart{Ch. 4 --- Disk, Gaussian and first-digit simulation}
Use independent $U,V\sim U(0,1)$ throughout.
\textbf{Uniform disk:} $R=\sqrt U$, $\Theta=2\pi V$; output $(R\cos\Theta,R\sin\Theta)$. Uniform radius overweights the centre because area grows as $r^2$.

\textbf{Box--Muller:} instead take $R=\sqrt{-2\ln U}$ with the same angle. The two coordinates are independent $N(0,1)$; rescale each by $\mu+\sigma Z$ as required.

\textbf{Benford:} when the fractional part of $\log_{10}X$ is uniform on $[0,1)$, first digit $D$ satisfies
\[
P(D=d)=\log_{10}(1+1/d),\qquad d=1,\ldots,9.
\]
Translate digit $d$ to fractional log in $[\log_{10}d,\log_{10}(d+1))$. This uniform-log condition is an assumption, not a rule for every dataset.
\boxend
\end{multicols*}

\newpage
% PAGE 2 / COLUMN 1: sampling and reference distributions.
\begin{multicols*}{3}
\vspace*{-1em}
\boxstart{ST2334 Finals --- Ch. 5--7: Statistics}
\textbf{Notation:} population parameters $\mu,\sigma^2$ are fixed; sample statistics $\bar X,S^2$ are random; observed values $\bar x,s^2$ are numbers. FS p.\,2 gives CI/test formulas for means. Use this page to choose and apply them. Here $q_\alpha$ always denotes an \emph{upper-tail} critical value.
\boxend

\boxstart{Ch. 5 --- Sampling, statistics and standard error}
\textbf{SRS without replacement:} every size-$n$ subset of a finite population is equally likely; equal individual inclusion probabilities alone are insufficient. Draws are dependent. \textbf{Iid sample:} independent draws from the same distribution, as with uniform sampling with replacement.

A statistic is a function of observations, with no unknown parameters. Its sampling distribution describes variation across repeated samples. Random sampling/randomisation supports inference; large sample size cannot remove selection bias.
\[
\bar X=\frac1n\sum_iX_i,\qquad
S^2=\frac{\sum_i(X_i-\bar X)^2}{n-1}
=\frac{\sum_iX_i^2-n\bar X^2}{n-1}.
\]
For iid observations with finite variance,
\[
E(\bar X)=\mu,\quad V(\bar X)=\sigma^2/n,\quad
\operatorname{SE}(\bar X)=\sigma/\sqrt n.
\]
Use $s/\sqrt n$ to \emph{estimate} the SE when $\sigma$ is unknown. SD measures spread of individual data; SE measures spread of an estimator. Halving SE requires four times as many observations.
\boxend

\boxstart{Ch. 5 --- LLN versus CLT; exact versus approximate}
\textbf{LLN:} $P(|\bar X_n-\mu|>\varepsilon)\to0$ for each $\varepsilon>0$: averages concentrate around $\mu$. It does not claim a normal shape.

\textbf{CLT:} for iid data with $0<\sigma^2<\infty$,
\[
\frac{\bar X-\mu}{\sigma/\sqrt n}\xrightarrow{d}N(0,1),\quad
\bar X\approx N(\mu,\sigma^2/n),\quad
\sum_iX_i\approx N(n\mu,n\sigma^2).
\]
Normal populations make the last two distributions \emph{exact} for every $n$. Otherwise the course's $n\ge30$ rule is approximate: skewness/heavy tails can require larger samples. CLT does not make the original observations normal, nor justify infinite-variance or dependent iid formulas.

\textbf{Chebyshev check:} $P(|\bar X-\mu|\ge\varepsilon)\le\sigma^2/(n\varepsilon^2)$ for iid finite-variance data, giving an LLN proof without normality.
\boxend

\boxstart{Ch. 5 --- CLT probabilities and sample-size questions}
\begin{enumerate}
\item Identify the target: one observation, a sample \emph{mean}, or a \emph{sum}.
\item Write its mean and variance; take the square root for its SD.
\item Standardise every bound; subtract CDF values for intervals.
\item For an unknown $n$, choose the correct tail quantile, solve the inequality, then round up and check the integer.
\end{enumerate}
\textbf{Example:} $\mu=120,\sigma=24,n=64$ gives SE $3$, so $P(\bar X>126)\approx P(Z>2)=0.0228$.

\textbf{One-sided design:} $\mu=80,\sigma=24$; require $P(\bar X>86)\le0.05$. Then $\sqrt n/4\ge z_{0.05}=1.645$, giving $n\ge43.2964$, hence $44$. Use $z_{\alpha}$ for one tail, $z_{\alpha/2}$ for central coverage.
\boxend

\boxstart{Ch. 5 --- Chi-squared and sample-variance events}
Independent $Z_i\sim N(0,1)$ give $U=\sum_{i=1}^\nu Z_i^2\sim\chi^2_\nu$; $E(U)=\nu$, $V(U)=2\nu$. Independent chi-squared variables add their degrees of freedom. Support is positive; the distribution is not symmetric.

For an iid \textbf{normal} sample,
\[
Q=\frac{(n-1)S^2}{\sigma^2}\sim\chi^2_{n-1};\qquad \bar X\text{ and }S^2\text{ are independent}.
\]
For $P(S^2>c)=\alpha$, set
$c=\sigma^2\chi^2_{n-1,\alpha}/(n-1)$.
For $P(S^2<c)=p$, replace the upper-tail index by $1-p$. For $S$ rather than $S^2$, square a positive cutoff first.

\textbf{Example:} $n=6,\sigma^2=4$; $P(S^2>c)=0.05$ gives $c=4(11.07)/5=8.856$. \\
\warn{$\frac{(n-1)S^2}{\sigma^2}\sim\chi^2_{n-1}$ needs normality; $E(S^2)=\sigma^2$ does not.}
\boxend

\boxstart{Ch. 5 --- Student's t and central probabilities}
For independent $Z\sim N(0,1)$ and $U\sim\chi^2_\nu$,
\[
T=Z/\sqrt{U/\nu}\sim t_\nu.
\]
The distribution is symmetric, has heavier tails than normal, and approaches normal as $\nu$ grows. For normal samples, replacing $\sigma$ by $S$ in the standardised mean gives exact $t_{n-1}$ (see FS).

\textbf{Transform an event:} if $n=25$, then $|\bar X-\mu|<cS$ becomes $|T|<5c$. For central probability $0.95$, $5c=t_{24,0.025}=2.064$, so $c=0.4128$.
\boxend

\boxstart{Ch. 5 --- F ratios: scaling, order and reciprocals}
For independent $U\sim\chi^2_a,V\sim\chi^2_b$, $F=(U/a)/(V/b)\sim F_{a,b}$. For independent normal samples,
\[
\frac{S_1^2/\sigma_1^2}{S_2^2/\sigma_2^2}\sim F_{n_1-1,n_2-1}.
\]
Thus $S_1^2/S_2^2>c$ becomes $F>c\sigma_2^2/\sigma_1^2$. It is an unscaled F ratio only when the population variances are equal. Numerator df comes first.
\[
1/F_{a,b}\sim F_{b,a},\qquad
F_{a,b;1-\alpha}=1/F_{b,a;\alpha}.
\]
Reciprocation reverses the inequality \emph{and} swaps df: $P(F_{8,15}>4/3)=P(F_{15,8}<3/4)$. Do not use symmetric-tail shortcuts for F or chi-squared.
\boxend

\columnbreak
% PAGE 2 / COLUMN 2: estimation and confidence intervals.
\boxstart{Ch. 6 --- Point estimates, bias and efficiency}
An \textbf{estimator} is a statistic; an \textbf{estimate} is its observed value. To test unbiasedness, compute $E_\theta(T)$ and compare with the \emph{target parameter}, not necessarily $\mu$. If $E(T)=a\theta$, $T/a$ is unbiased when $a\ne0$ is known.

Examples: for $U(0,\theta)$, $\bar X$ estimates $\theta/2$, so $2\bar X$ is unbiased for $\theta$. For Binomial$(m,p)$ data with known $m$, $\bar X/m$ is unbiased for $p$. Among unbiased estimators of the same target, compare variances; lower is better. Unbiasedness does not imply small sampling error.
\boxend

\boxstart{Ch. 5/6 --- Why the sample variance uses $n-1$}
For iid data with mean $\mu$ and finite variance $\sigma^2$,
\[
\sum_i(X_i-\bar X)^2=\sum_i(X_i-\mu)^2-n(\bar X-\mu)^2.
\]
Take expectations: $n\sigma^2-n(\sigma^2/n)=(n-1)\sigma^2$, hence $E(S^2)=\sigma^2$. Dividing by $n$ instead produces bias factor $(n-1)/n$.

This proof needs no normality. $S^2$ unbiased for $\sigma^2$ does \emph{not} make $S$ unbiased for $\sigma$: expectations do not generally commute with square roots.
\boxend

\boxstart{Ch. 6 --- Uniform extrema: tutorial estimator toolkit}
For iid $U(0,b)$ data, let $L=\min X_i$, $M=\max X_i$. For $0<x<b$,
\[
F_M(x)=(x/b)^n,\quad P(L>x)=(1-x/b)^n,
\]
\[
E(M)=\frac{nb}{n+1},\quad E(L)=\frac{b}{n+1},\quad
V(M)=V(L)=\frac{nb^2}{(n+1)^2(n+2)}.
\]
Derive densities by differentiation and integrate for moments. For $b=\theta$, $(n+1)M/n$ is unbiased with variance $\theta^2/[n(n+2)]$; $2\bar X$ is unbiased with variance $\theta^2/(3n)$. The maximum-based estimator has smaller variance for $n>1$.

\textbf{Linear combinations (Tutorial 9):} two iid $U(0,2\theta)$ observations have $\operatorname{Cov}(L,M)=\theta^2/9$. For $T=aL+bM$ (here $a,b$ are weights),
\[
E(T)=\tfrac{2\theta}{3}(a+2b),\quad
V(T)=\tfrac{2\theta^2}{9}(a^2+ab+b^2).
\]
Unbiasedness imposes $a+2b=3/2$; substitute into the variance and minimise the quadratic. The optimum is $a=0,b=3/4$, variance $\theta^2/8$. \\
\warn{Extrema of one sample are dependent.}
\boxend

\boxstart{Ch. 6 --- Precision and minimum sample size}
For desired two-sided margin $e>0$ with known/planning SD $\sigma$ and normal/large-sample model,
\[
n_{\min}=\left\lceil\left(\frac{z_{\alpha/2}\sigma}{e}\right)^2\right\rceil.
\]
If given total interval width $w$, use $e=w/2$. Increasing confidence widens the interval; increasing $n$ narrows it as $1/\sqrt n$; halving the margin requires four times $n$.

Unknown $\sigma$ requires a prior/pilot estimate for planning. For a small-sample t margin, the critical value itself depends on $n$: solve iteratively and check the chosen df. FS p.\,2 supplies the margin as the CI half-width.

\textbf{Example:} 99\% confidence, $\sigma=3.5$, $e=0.5$: $n\ge(2.576\cdot3.5/0.5)^2=325.15$, so choose $326$.
\boxend

\boxstart{Ch. 6 --- Confidence intervals: use FS p. 2}
\begin{enumerate}
\item Identify the target: one mean, independent mean difference, or mean of paired differences.
\item Check the design, known versus estimated population variances, normality, and each sample size.
\item Choose the matching FS row and its standard error; find the upper-tail $\alpha/2$ critical value and df if needed.
\item Compute endpoints, preserving the subtraction order and measurement units; interpret in context.
\end{enumerate}
\textbf{Course convention:} large means $n\ge30$; \emph{both} sizes must qualify for independent large-sample formulas. Unknown $\sigma$ with small normal samples uses t. Small non-normal samples are not covered.

\textbf{Interpretation:} $100(1-\alpha)\%$ of repeated-sample intervals cover the fixed parameter. This is not a post-data probability about that parameter, nor coverage of individual observations.
\boxend

\boxstart{Ch. 6/7 --- Independent means: assumptions first}
Use FS Section 2, preserving the order $\mu_1-\mu_2$. Known variances with normal populations give exact normal inference; large non-normal samples give an approximation.

\textbf{Unknown variances:} large samples allow the unpooled normal approximation. Equal population variances justify pooling, weighted by each sample's df; use the \emph{square root} of the pooled variance for its SD. Small normal samples with equal variances use pooled t, df $n_1+n_2-2$.

Sample SDs alone do not prove/refute equal population variances. Pool only when justified; FS has no small-sample unequal-variance row.
\boxend

\boxstart{Ch. 6/7 --- Paired data: analyse the differences}
Same unit before/after or matched units $\Rightarrow$ define $D_i=X_i-Y_i$. Use FS Section 1 on this \emph{single} sample: $n$ counts pairs; small-sample df is $n-1$.

Compute $s_D$ from the differences, not from the two marginal SDs. Within-pair covariance matters: $V(D)=V(X)+V(Y)-2\operatorname{Cov}(X,Y)$. Different pairs must be independent; for small samples, normality concerns the \emph{differences}.

For $D=$ after $-$ before, improvement means $\mu_D>0$ for scores but $\mu_D<0$ for latency. Reversing $D$ reverses the estimate and alternative; CI $(L,U)$ becomes $(-U,-L)$.
\boxend

\columnbreak
% PAGE 2 / COLUMN 3: hypothesis testing and exam application.
\boxstart{Ch. 7 --- Hypothesis-test answer template}
\begin{enumerate}
\item Define the population parameter; state $H_0$ and directional or two-sided $H_1$ \emph{before} examining the result.
\item State $\alpha$, the assumptions and the chosen FS statistic.
\item Specify its distribution under $H_0$ and df where needed.
\item Calculate the observed statistic; obtain a p-value or compare with the appropriate critical value.
\item Reject or fail to reject $H_0$; conclude about the population in the question's wording and units.
\end{enumerate}
For a nonzero hypothesised difference $\delta_0$, subtract it from $\bar x-\bar y$ in the FS numerator. For paired tests use $\bar d-\delta_0$. ``Changed/different'' is two-sided; ``increased/greater'' and ``decreased/less'' specify one tail.
\boxend

\boxstart{Ch. 7 --- Rejection regions and p-values}
Let $W$ have the continuous symmetric null distribution (standard normal or t), $G$ its CDF, $w$ the observed statistic, and $q_\gamma$ its upper-tail $\gamma$ critical value.
\[
\begin{array}{c|c|c}
H_1&\text{reject when}&p\text{-value}\\\hline
\theta>\theta_0&w>q_\alpha&1-G(w)\\
\theta<\theta_0&w<-q_\alpha&G(w)\\
\theta\ne\theta_0&|w|>q_{\alpha/2}&2[1-G(|w|)]
\end{array}
\]
Use $z_\gamma$ for normal, $t_{\nu,\gamma}$ for t. The course rule is reject when $p<\alpha$. A statistic in the opposite direction to a one-sided alternative gives a large p-value; do not take its absolute value.

\textbf{Definition:} probability \emph{under $H_0$} of a result at least as extreme in the alternative's direction(s). It is not $P(H_0\mid\text{data})$, the effect size, or the probability of a Type I error after seeing the data.
\boxend

\boxstart{Ch. 5/7 --- Reading the supplied critical-value tables}
\begin{itemize}
\item Normal CDF entries are left-tail areas. Critical-value tables for t, F and chi-squared use right-tail areas: read each heading.
\item For two-sided normal/t inference, put $\alpha/2$ in \emph{each} tail. For t, select df, not sample size.
\item If $q_{0.05}<w<q_{0.025}$ for positive $w$, the upper-tail p-value lies between $0.025$ and $0.05$; a two-sided symmetric p-value lies between $0.05$ and $0.10$.
\item F tables: numerator df across columns, denominator df down rows. Reverse both dfs for a reciprocal.
\item For untabulated values, report a supported bound or distribution expression. The normal limit approximates t at finite df.
\end{itemize}
\boxend

\boxstart{Ch. 7 --- Errors, significance and power}
\begin{tabularx}{\linewidth}{@{}lXX@{}}
Reality & Fail to reject $H_0$ & Reject $H_0$\\\hline
$H_0$ true & Correct & Type I error ($\alpha$)\\
$H_0$ false & Type II error ($\beta$) & Correct (power $1-\beta$)
\end{tabularx}
\textbf{Power is evaluated at a specified true alternative.} Build the rejection region under $H_0$, then compute its probability under that alternative. For a known-SD right-tailed mean test, let $s_e=\sigma/\sqrt n$ and $c=\mu_0+z_\alpha s_e$. At true mean $\mu_1$,
\[
\beta(\mu_1)=\Phi\!\left(\frac{c-\mu_1}{s_e}\right),\qquad
\text{power}=1-\beta(\mu_1).
\]
For a two-sided test with acceptance interval $[L,U]$ for $\bar X$, $\beta(\mu_1)=\Phi((U-\mu_1)/s_e)-\Phi((L-\mu_1)/s_e)$. These normal calculations are exact for normal data and approximate under CLT.

At fixed effect and $n$, smaller $\alpha$ generally increases $\beta$. Larger $n$, larger effect or smaller noise generally increases power. \\
\warn{Failing to reject is not proof of equality.}
\boxend

\boxstart{Ch. 6/7 --- Link confidence intervals to tests}
A two-sided level-$\alpha$ test rejects exactly when the hypothesised value lies outside the matching $100(1-\alpha)\%$ CI, using the same model and SE. For differences compare with $\delta_0$ (usually 0).

For $\mu_1-\mu_2$, intervals wholly above/below 0 support positive/negative differences; containing 0 does not establish equivalence. The rule does not apply unchanged to one-sided tests. Significance need not imply a practically large effect.
\boxend

\boxstart{Ch. 6/7 --- Compact worked decisions}
\textbf{One mean, small normal sample:} $n=7$, $\bar x=10$, $s^2=0.08$. Unknown variance $\Rightarrow$ t, df $6$; for 95\% CI use $t_{6,0.025}=2.447$:
\[
10\pm2.447\sqrt{0.08/7}=(9.738,10.262).
\]
\textbf{Independent, small, normal, equal variances:} $n_1=12,n_2=10$, means $85,81$, SDs $4,5$. FS pooling gives $s_p^2=20.05$, df $20$; 90\% CI:
\[
4\pm1.7247\sqrt{20.05(1/12+1/10)}=(0.693,7.307).
\]
\textbf{Paired improvement:} $n=8$, $\bar d=24.25$, $s_D=25.39$, normal differences (after $-$ before). Test $H_0:\mu_D=0$ vs $H_1:\mu_D>0$: $t=24.25/(25.39/\sqrt8)=2.70$, df $7$. Since $2.70>t_{7,0.05}=1.895$, reject at 5\%: evidence of increased population mean.
\boxend

\boxstart{All chapters --- Final answer checks}
\textbf{Probability:} support, endpoints, conditioning denominator, units, range $[0,1]$, and independent versus disjoint events.
\textbf{Statistics:} population/sample SD; mean/sum SE; paired/independent; known/unknown variance; df and tail; subtract the null value; round $n$ up; conclude about the population.
\boxend
\end{multicols*}
\end{document}
